\section*{Kapitel 5 --- Framen i kod, och den första testsviten}
\addcontentsline{toc}{section}{Kapitel 5 --- Framen i kod}

\solution{ex:5:1}{Förutsäg}
\ans{a} Typiskt 12 bytes, inte 11. Kompilatorn lägger in utfyllnad så att \code{id}
(\code{std::uint16_t}) hamnar på jämn adress, och avslutar strukten på en storlek som är en
multipel av dess strängaste alignment.
\ans{b} Ja. \code{f.data[7]} är \code{0}: aggregate-initiering nollställer de medlemmar som inte
anges.
\ans{c} Raden slutar kompilera. En användardeklarerad konstruktor gör typen till en icke-aggregate,
och då finns ingen konstruktor som tar \code{(0x123U, 2U, {...})}. Det är precis skälet till att
\code{Frame} inte har någon.

\solution{ex:5:2}{constexpr}
\ans{a} Att \code{isValid()} går att utvärdera vid kompileringstillfället --- annars är den inte
ett giltigt \code{static_assert}-villkor. Det bevisar samtidigt att funktionen är fri från
sidoeffekter.
\ans{b} Ingenting. Kontrollen sker när programmet byggs; i den färdiga binären finns ingen kod för
den alls.

\solution{ex:5:3}{Läs sviten}
\ans{a} \code{0x7FF} och \code{0x800} för \code{id}, \code{8} och \code{9} för \code{dlc} --- det
största som ska accepteras och det minsta som ska avvisas. Buggar bor på gränserna: ett värde mitt
i intervallet passerar både en korrekt och en felaktig jämförelse.
\ans{b} Fallet som kontrollerar att \code{dlc = 8} accepteras skulle fallera; alla andra vore
gröna. Det är därför paret behövs och inte bara det ena värdet.
\ans{c} Ja --- att \code{MaxDataLen} är samma tal som arrayens storlek kontrolleras inte, och
kan inte kontrolleras utifrån. En grön svit betyder att koden är rätt på de punkter sviten frågar
om, inte att den är rätt.

\solution{ex:5:4}{Bryt koden}
\ans{a} Fallet som accepterar det största giltiga \code{id}:t, alltså \code{0x7FF}. Utskriften
namnger uttrycket, filen och raden.
\ans{b} Kortare än att gissa: testnamnet pekar ut kravet, Arrange-delen förutsättningen, och
Assert-raden det förväntade värdet --- först därefter behöver man öppna sin egen kod.
\ans{c} Återställ koden --- och notera att den felaktiga versionen passerade alla andra testfall.
Ett enda tecken skilde en korrekt implementation från en som avvisar var åttonde giltig frame.

\solution{ex:5:5}{Var ligger valideringen}
\ans{a} Att ett fel fångas tidigare och närmare den som orsakade det.
\ans{b} Att regeln nu finns på två ställen, och att båda måste hållas i synk.
\ans{c} Bara det ena stället ändras. Den glömda kopian avvisar då frames som är giltiga, eller
släpper igenom frames som inte är det --- och felet ser ut att komma från hårdvaran.

Regeln uttrycks en gång, i \code{isValid()}, och används på ett ställe: i \code{Spi::send()}.
